% unit: noether.p12.title.001
\providecommand{\bdelta}{\bar{\delta}}
\begin{center}
{\Large\bfseries 12. Invariants d'expressions différentielles arbitraires}\par
\vspace{0.8em}
Par Emmy Noether, à Göttingen.\par
\vspace{1.5em}
\emph{Nachrichten von der Gesellschaft der Wissenschaften zu Göttingen} 1918, pp. 37--44
\end{center}

\vspace{3em}
\begin{center}
Présenté dans la séance du 25 janvier 1918 par F. Klein.
\end{center}

% unit: noether.p12.par.001
Par expression différentielle j'entends une fonction
% unit: noether.p12.eq.000
\[
f(x,dx)=f(x_1,\ldots,x_n;\,dx_1,\ldots,dx_n)
\]
des \(n\) variables \(x_1,\ldots,x_n\) et de leurs différentielles
\(dx_1,\ldots,dx_n\), que l'on suppose analytique dans les deux systèmes
de \(n\) arguments chacun, mais non nécessairement réelle. Je prends
maintenant pour base le groupe de toutes les transformations analytiques
des variables, auquel correspond le groupe de toutes les transformations
linéaires des différentielles:
% unit: noether.p12.eq.001
\begin{equation}
x_i=x_i(y_1,\ldots,y_n);\qquad
dx_i=\sum_k\frac{\partial x_i}{\partial y_k}\,dy_k;\qquad
\delta x_i=\sum_k\frac{\partial x_i}{\partial y_k}\,\delta y_k,
\tag{1}
\end{equation}
par lesquelles \(f(x,dx)\) doit passer en \(g(y,dy)\). Par invariant de
\(f(x,dx)\), j'entends un invariant relativement à ce groupe, et relativement
au groupe ainsi induit pour les différentielles supérieures
\(d^2x,d\delta x,\ldots\) et pour les dérivées de \(f(x,dx)\); donc une fonction
\(J\) (analytique) telle que, pour \(f\) arbitraire, et en vertu de (1),
on ait identiquement, dans les variables et dans toutes les différentielles
qui interviennent:
% unit: noether.p12.eq.002
\[
\begin{aligned}
&J\!\left(f,\frac{\partial f}{\partial dx}\cdots
\frac{\partial^{\rho+\sigma}f}{\partial x^\rho\partial dx^\sigma}
\cdots dx,\delta x,d^2x,\ldots\right)\\
&\qquad =
J\!\left(g,\frac{\partial g}{\partial dy}\cdots
\frac{\partial^{\rho+\sigma}g}{\partial y^\rho\partial dy^\sigma}
\cdots dy,\delta y,d^2y,\ldots\right).\footnotemark
\end{aligned}
\]
\footnotetext{\(\displaystyle
\frac{\partial^{\rho+\sigma}f}{\partial x^\rho\partial dx^\sigma}\)
sera partout employé comme abréviation de
\[
\frac{\partial^{\rho+\sigma}f}
{\partial x_{i_1}\cdots\partial x_{i_\rho}
 \partial dx_{k_1}\cdots\partial dx_{k_\sigma}},
\]
où les indices désignent quelconques des nombres \(1\) à \(n\).}

% unit: noether.p12.par.002
On définit exactement de même l'invariant d'un système simultané
d'expressions différentielles. Si, en particulier, un tel invariant ne contient
que les premières différentielles \(dx,\delta x\), et seulement des dérivées
par rapport à ces différentielles, non par rapport aux variables, alors
l'apparition des variables dans les fonctions et la transformation linéaire
des différentielles n'entrent pas en considération; ces invariants spéciaux
sont des invariants relativement au groupe de toutes les transformations
linéaires des différentielles \(dx,\delta x\), à coefficients indéterminés,
et seront appelés invariants projectifs du système simultané.

% unit: noether.p12.par.003
Pour les formes différentielles homogènes quadratiques, Christoffel
(Crelle 70) et Ricci (\emph{Math. Annalen} 54) ont établi un système de
formations invariantes tel que tous les invariants deviennent des invariants
projectifs de ce système simultané; par là, les questions relatives à la
totalité des invariants et à l'équivalence se ramènent à des questions de
théorie linéaire des invariants. C'est ce que je voudrais appeler le
\og{} théorème de réduction \fg{}. Le système complet se compose d'une infinité
dénombrable de formes; mais, pour les invariants d'un ordre fini quelconque
\(\rho\), c'est-à-dire ceux qui ne contiennent pas de dérivées par rapport
aux variables d'ordre supérieur au \(\rho\)-ième, il ne se compose que d'un
nombre fini de formes.

% unit: noether.p12.par.004
Dans ce qui suit, je montre la validité du \og{} théorème de réduction \fg{} pour
des expressions différentielles arbitraires; ici encore il s'agit d'un système
complet d'une infinité dénombrable de formations invariantes, mais, pour
chaque ordre fini, seulement d'un nombre fini d'entre elles. Toutefois, alors
que Christoffel et Ricci fondent la formation des invariants et la preuve du
théorème de réduction sur des éliminations difficiles à embrasser d'un coup
d'\oe{}il, je forme les invariants directement par des procédés de différentiation
invariante et de variation, ces derniers pouvant être regardés simplement
comme des procédés de différentiation par rapport à d'autres paramètres.
L'algorithme de ces procédés de différentiation invariante fut développé,
à la suite de Lagrange (\emph{Mécanique analytique}, deuxième partie,
section IV), par Riemann (\OE{}uvres XXII) et Lipschitz (Crelle 70, 72) pour
les formes différentielles quadratiques, sans toutefois aller jusqu'au
théorème de réduction. Les moyens auxiliaires pour prouver le théorème de
réduction sont fournis par les \og{} coordonnées normales \fg{} introduites par
Riemann pour les formes quadratiques (\OE{}uvres XIII); elles transforment
en droites les extrémales, issues d'un point, du problème variationnel
associé, mais elles n'existent que pour les expressions différentielles
homogènes
\(f(x,\varkappa\cdot dx)=\Phi(\varkappa)\cdot f(x,dx)\footnote{Il se voit
facilement que \(\Phi(\varkappa)=\varkappa^c\), où \(c\) désigne une
grandeur réelle ou complexe quelconque.}\). Le cas inhomogène se ramène
cependant à celui-ci.

% unit: noether.p12.heading.001
\begin{center}
\textbf{I. Formation des invariants.}
\end{center}

% unit: noether.p12.par.005
Une suite d'invariants projectifs de \(f(x,dx)\), qui en général ne s'arrête
pas, est fournie par les polaires. À cause de la linéarité de la transformation
des différentielles, de \(f(x,dx)=g(y,dy)\), on déduit aussi que
\(f(x,dx+\lambda\delta x)=g(y,dy+\lambda\delta y)\); et de là, par
développement suivant \(\lambda\), naissent les relations caractéristiques
de l'invariance:
% unit: noether.p12.eq.003
\begin{equation}
\begin{gathered}
f(dx)=g(dy),\\
f_\delta=\sum\frac{\partial f(dx)}{\partial dx}\,\delta x
       =\sum\frac{\partial g(dy)}{\partial dy}\,\delta y,\\
\vdots\\
f_{\delta^\sigma}=
\sum\frac{\partial^\sigma f(dx)}{\partial dx^\sigma}\,\delta x^\sigma
=
\sum\frac{\partial^\sigma g(dy)}{\partial dy^\sigma}\,\delta y^\sigma,\\
\vdots
\end{gathered}
\tag{2}
\end{equation}

% unit: noether.p12.par.006
Chaque polaire, par application d'un procédé de variation \(\bdelta\),
donne naissance à une suite non terminante d'invariants généraux:
% unit: noether.p12.eq.004
\begin{equation}
\begin{gathered}
\bdelta^\rho f(x,dx)=\bdelta^\rho g(y,dy),\\
\vdots\\
\bdelta^\rho f_{\delta^\sigma}
=\bdelta^\rho\sum\frac{\partial^\sigma f(dx)}{\partial dx^\sigma}\,
\delta x^\sigma
=\bdelta^\rho\sum\frac{\partial^\sigma g(dy)}{\partial dy^\sigma}\,
\delta y^\sigma,\\
\vdots\\[0.2em]
(\rho=0,1,2,\ldots).
\end{gathered}
\tag{3}
\end{equation}
Au moyen de (1), les relations (2) et (3) permettent de lire les lois de
transformation des dérivées de \(f(x,dx)\); cela ne sera cependant pas
nécessaire dans la suite. À partir des invariants (3), on peut maintenant
obtenir par combinaison linéaire la \og{} forme normale de la \(\rho\)-ième
variation \fg{}, que l'on trouve pour \(\rho=1\) chez Lagrange et pour
\(\rho=2\) chez Riemann; elle est caractérisée par le fait que les
différentielles mixtes d'ordre \(\rho+1\), à savoir
\(d^\rho\delta,d^{\rho-1}\delta^2,\ldots\), sont éliminées. On obtient:
% unit: noether.p12.eq.005
\begin{equation}
\begin{gathered}
\Omega_1=\delta f(dx)-d f_\delta(dx),\\
\Omega_2=\delta^2 f-\delta d f_\delta+\frac12 d^2 f_{\delta^2},\\
\vdots\\
\Omega_\rho=\delta^\rho f-\delta^{\rho-1}d f_\delta
+\frac12\delta^{\rho-2}d^2 f_{\delta^2}
+\cdots+\frac{(-1)^\rho}{\rho!}\,d^\rho f_{\delta^\rho},\\
\vdots
\end{gathered}
\tag{4}
\end{equation}

% unit: noether.p12.par.007
À partir des invariants (4), en généralisant un ansatz de Riemann, on peut
maintenant parvenir à des invariants qui ne contiennent que des premières
différentielles; de tels invariants seront appelés \og{} fonctions fondamentales \fg{}.
Si, en effet, dans \(\Omega_1\), les différentielles sont remplacées par les
quotients différentiels, alors \(\Omega_1=0\) (identiquement en \(\bdelta x\))
donne exactement les équations de Lagrange du problème variationnel
appartenant à
% unit: noether.p12.eq.006
\[
f\!\left(\frac{dx}{dt}\right).
\]
J'impose maintenant la condition invariante que la direction
\((d+\lambda\delta)\) satisfasse à ces équations de Lagrange:
% unit: noether.p12.eq.007
\begin{equation}
\Omega_1(d+\lambda\delta)=
\bdelta f(dx+\lambda\delta x)
-(d+\lambda\delta)f_{\bdelta}(dx+\lambda\delta x)=0
\quad(\text{identiquement en }\bdelta x),
\tag{5}
\end{equation}
d'où, par la substitution \(\bdelta=d+\lambda\delta\), il résulte encore,
pour \(f(dx)\) homogène:
% unit: noether.p12.eq.008
\begin{equation}
(d+\lambda\delta)f(dx+\lambda\delta x)=0,
\tag{6}
\end{equation}
sauf dans le cas de l'homogénéité du premier ordre. Dans ce dernier cas,
(6) doit être ajoutée comme nouvelle condition, car alors le système
d'équations (5) ne représente que \(n-1\) conditions indépendantes. Si,
dans (5), respectivement dans (5) et (6), on annule les coefficients des
puissances zéro, première et seconde de \(\lambda\), on obtient ainsi
trois systèmes invariants d'équations
\(\varphi=0\), \(\psi=0\), \(\chi=0\) (identiquement en \(\bdelta x\)),
au moyen desquels \(d^2x,d\delta x,\delta^2x\) peuvent être exprimés par les
premières différentielles. Les autres systèmes invariants d'équations
% unit: noether.p12.eq.009
\begin{equation}
\begin{gathered}
d^\sigma\varphi=0,\qquad d^\sigma\psi=0,\qquad
d^\sigma\chi=0,\qquad d^{\sigma-1}\delta\chi=0,\\
d^{\sigma-2}\delta^2\chi=0,\ldots,\delta^\sigma\chi=0
\qquad(\sigma=0,1,2,\ldots)
\end{gathered}
\tag{7}
\end{equation}
suffisent alors précisément pour exprimer toutes les différentielles
supérieures par les premières.\footnote{Seule la forme linéaire
\(\sum A_i\,dx_i\) demande un traitement particulier, car ici le système
d'équations (5) ne contient pas de secondes différentielles.} Les fonctions
(4), liées au système invariant d'équations (7), conduisent donc à des
fonctions fondamentales \([\Omega_\rho]\) ne contenant que des premières
différentielles; de plus \([\Omega_1]\) s'annule identiquement.

% unit: noether.p12.par.008
De même, à partir de la différentielle \(\delta h(dx,\delta x)\) d'une
fonction fondamentale \(h\), on peut, au moyen de (7), obtenir de nouveau
une fonction fondamentale, qui sera appelée la \og{} dérivée covariante \fg{}
\([h^{(1)}]\) de \(h\); en général, \([h^{(\nu)}]\) désignera la dérivée
covariante de \([h^{(\nu-1)}]\). Ces deux procédés conduisent ainsi à
une suite doublement infinie de fonctions fondamentales:
% unit: noether.p12.eq.010
\begin{equation}
\begin{array}{cccccc}
[\Omega_2], & [\Omega_2^{(1)}], & [\Omega_2^{(2)}], & \ldots,
& [\Omega_2^{(\nu)}], & \ldots\\
\vdots &&&& \vdots&\\
{}[\Omega_\rho], & [\Omega_\rho^{(1)}], & \ldots, & \ldots,
& [\Omega_\rho^{(\nu)}], & \ldots\\
\vdots & \vdots &&& \vdots&
\end{array}
\tag{8}
\end{equation}

% unit: noether.p12.par.009
Il apparaît en outre que les membres de gauche des équations qui expriment
les différentielles supérieures par les premières sont cogrédients aux
premières différentielles, c'est-à-dire qu'ils subissent les mêmes
transformations linéaires. Si donc, à la place des différentielles supérieures,
on introduit dans l'invariant ces membres de gauche \(p,q,r,\ldots\)
comme arguments, alors, abstraction faite des dérivées, il ne dépend que
de grandeurs cogrédientes entre elles. La formation des fonctions (8)
indiquée ci-dessus revient simplement à l'introduction de ces nouvelles
grandeurs; chaque invariant passe ainsi en une somme d'invariants, et de
cette somme on choisit celui qui contient précisément \(dx,\delta x\), mais
non \(p,q,r,\ldots\).

% unit: noether.p12.par.010
Dans la suite on montre que, sous l'hypothèse d'homogénéité
% unit: noether.p12.eq.011
\[
f(\varkappa\cdot dx)=\Phi(\varkappa)\cdot f(dx),
\]
les fonctions fondamentales (8) épuisent le système complet cherché.

% unit: noether.p12.heading.002
\begin{center}
\textbf{II. Coordonnées normales, équivalence et théorème de réduction.}
\end{center}

% unit: noether.p12.par.011
J'arrive aux coordonnées normales en intégrant le problème variationnel
appartenant à \(f\!\left(\frac{dx}{dt}\right)\):
% unit: noether.p12.eq.012
\begin{equation}
\begin{gathered}
\Omega_1=\delta f\!\left(\frac{dx}{dt}\right)
-\frac{d}{dt}f_\delta\!\left(\frac{dx}{dt}\right)=0
\quad(\text{identiquement en }\delta x),\\
\frac{d}{dt}f\!\left(\frac{dx}{dt}\right)=0
\quad(\text{comme conséquence ou condition additionnelle}).
\end{gathered}
\tag{9}
\end{equation}
En raison de l'homogénéité supposée de \(f(dx)\), ce système d'équations
se transforme en lui-même par la substitution de paramètre
\(t=\varkappa\cdot\tau\). Par conséquent, si, au moyen de (9), on exprime
\(\frac{d^\rho x_i}{dt^\rho}\) comme fonction
\(\varphi_\rho^{(i)}\!\left(\frac{dx}{dt}\right)\) des premières dérivées,
alors \(\varphi_\rho^{(i)}\) est homogène d'ordre \(\rho\):
% unit: noether.p12.eq.013
\begin{equation}
\varphi_\rho^{(i)}\!\left(\varkappa\cdot\frac{dx}{dt}\right)
=\varkappa^\rho\varphi_\rho^{(i)}\!\left(\frac{dx}{dt}\right).
\tag{10}
\end{equation}
Si, dans l'intégration de (9), on prescrit maintenant, pour \(t=t_0\),
les valeurs initiales
% unit: noether.p12.eq.014
\[
x=(x)_0,\qquad \frac{dx}{dt}=\left(\frac{dx}{dt}\right)_0
\]
et si l'on pose
% unit: noether.p12.eq.015
\begin{equation}
u_i=(t-t_0)\left(\frac{dx_i}{dt}\right)_0,
\tag{11}
\end{equation}
où, en vertu de \(f\!\left(\frac{dx}{dt}\right)=\mathrm{const.}\), le
paramètre \((t-t_0)\) devient proportionnel à la \og{} longueur de l'extrémale \fg{},
alors le développement en puissances de \((t-t_0)\), en vertu de (10),
donne
% unit: noether.p12.eq.016
\begin{equation}
x_i-(x_i)_0=u_i+\frac12\varphi_2^{(i)}(u)+\cdots+
\frac{1}{\rho!}\varphi_\rho^{(i)}(u)+\cdots .
\tag{12}
\end{equation}
Ce développement peut être regardé comme une transformation de variables
\(x_i=x_i(u)\), où les grandeurs \(u\) sont précisément les coordonnées
normales de Riemann. Ainsi, par (11) et (12), ces coordonnées transforment
les extrémales passant par \((x)_0\) en droites. De plus, à une transformation
arbitraire de variables \(x=x(y)\) correspond, par (11), une transformation
linéaire \(u=u(v)\) des coordonnées normales associées, dont les coefficients
ne dépendent que du système arbitraire de valeurs \((x)_0\); les différentielles
\(du,\delta u\) sont donc soumises à la même transformation que les \(u\)
eux-mêmes.

% unit: noether.p12.par.012
Au moyen de (12), que \(f(x,dx)\) passe en \(F(u,du)\), ou, développé en
puissances des \(u\),
% unit: noether.p12.eq.017
\begin{equation}
F(u,du)=F_0(u,du)+F_1(u,du)+\cdots+F_\rho(u,du)+\cdots ,
\tag{13}
\end{equation}
où \(F_\rho(u,du)\) désigne une forme homogène de dimension \(\rho\) en
\(u\). À cause de la linéarité de la transformation \(u=u(v)\), l'identité
\(F(u,du)=G(v,dv)\) transporte également les composantes homogènes
particulières dans les composantes correspondantes:
% unit: noether.p12.eq.018
\begin{equation}
F_\rho(u,du)=G_\rho(v,dv).
\tag{14}
\end{equation}
Si donc on prend pour \((x)_0\) un système fixe quelconque de valeurs
constantes, (13) et (14) montrent que l'équivalence de \(f(dx)\) avec
\(g(dy)\) équivaut à l'équivalence des systèmes de fonctions associés
\(F_\rho(u,du)\) relativement à une transformation linéaire. En outre,
les \(F_\rho(u,du)\), ou les \(F_\rho(\delta u,du)\) correspondantes,
représentent des invariants de \(f(dx)\), pris au point \(x=(x)_0\);
elles forment même un système complet de fonctions fondamentales pour
ce point \(x=(x)_0\). En effet,
% unit: noether.p12.eq.019
\[
F_\rho(\delta u,du)=\frac1{\rho!}\sum
\frac{\partial F_\rho}{\partial \delta u^\rho}\,\delta u^\rho;
\qquad
\frac{\partial F_\rho}{\partial \delta u^\rho}
=\left(\frac{\partial F(du)}{\partial u^\rho}\right)_0,
\]
et cette dernière dérivée peut, au moyen de (12), s'exprimer de la même
manière par les dérivées de \(f(dx)\), prises pour \(x=(x)_0\), que la
dérivée correspondante de \(G(dv)\) peut s'exprimer par les dérivées de
\(g(dy)\). Comme
% unit: noether.p12.eq.020
\[
\left(\frac{\partial^{\rho+\sigma}F(du)}
{\partial u^\rho\partial du^\sigma}\right)_0
=
\frac{\partial^{\rho+\sigma}F_\rho(\delta u,du)}
{\partial\delta u^\rho\partial du^\sigma},
\]
chaque invariant de \(f(dx)=F(du)\), pris au point \(x=(x)_0\), devient
aussi un invariant projectif des \(F_\rho(\delta u,du)\), dès que les
différentielles supérieures ont seulement été remplacées par les
\(p,q,r,\ldots\) cogrédients à \(dx,\delta x\).

% unit: noether.p12.par.013
Mais maintenant \((x)_0\) peut être regardé comme un paramètre; à chaque
invariant au point \(x=(x)_0\) correspond un invariant de \(f(dx)\), pourvu
seulement que \((x)_0\) soit de nouveau remplacé par \(x\). Ainsi, si les
\(F_\rho(\delta u,du)\) deviennent des invariants
\(\Psi_\rho(\delta x,dx)\), pris pour \(x=(x)_0\) --- car, pour \(x=(x)_0\),
\(du,\delta u\) passent en \(dx,\delta x\) ---, alors les \(\Psi_\rho\)
elles-mêmes forment les fonctions fondamentales d'un système complet.
Il reste seulement à exprimer ce système de fonctions fondamentales par
des invariants explicites de \(f(dx)\), à savoir par les fonctions (8).
Que cela soit possible résulte de leur formation au moyen de procédés de
différentiation. Pour les termes d'ordre le plus élevé, ces procédés se
transforment en simples procédés de polarisation; et une transformation
analogue au développement en série de Clebsch--Gordan, combinée avec une
induction à cause des termes négligés, démontre l'assertion. Si l'on a
affaire aux invariants d'un système simultané, il faut encore ajouter aux
fonctions fondamentales (8) les dérivées covariantes des autres expressions
différentielles, comme le montre l'introduction, dans ces expressions, des
coordonnées normales appartenant à \(f(dx)\).

% unit: noether.p12.par.014
L'équivalence de \(f(dx)\) avec \(g(dy)\), qui avait été réduite à (14),
est donc aussi équivalente à l'équivalence des \(\Psi_\rho\), ou tout aussi
bien des fonctions (8), relativement à la transformation linéaire des
différentielles, sans qu'aucune condition spéciale d'intégrabilité soit
requise; cela résulte de la possibilité de la réduction à (14). Ainsi le
théorème de réduction est démontré dans toutes ses parties. En particulier,
l'annulation identique de la suite de fonctions
\([\Omega_2],[\Omega_3],\ldots,[\Omega_\rho],\ldots\) est nécessaire et
suffisante pour que \(f(dx)\) puisse être transformée en une expression à
coefficients constants.

% unit: noether.p12.par.015
Enfin, si l'on prend pour base, au lieu du groupe de toutes les transformations
analytiques, un sous-groupe, alors le groupe des transformations linéaires
correspondantes des \(u\) passe lui aussi en un sous-groupe du groupe
projectif; les invariants de \(f(dx)\) redeviennent des invariants des
fonctions (8) relativement à une transformation linéaire, mais maintenant
relativement à ce sous-groupe. Ainsi, par homogénéisation, le cas des
fonctions non homogènes \(f(dx)\) peut être ramené au groupe affine; le
système complet peut alors être déduit des fonctions (8) formées pour
une variable de plus.

% unit: noether.p12.par.016
Du théorème de réduction on obtient le cas particulier des formes quadratiques,
et plus généralement des formes de dimension \(p\), du fait qu'ici le système
de fonctions s'arrête à \([\Omega_2]\), plus généralement à \([\Omega_p]\)
et à ses dérivées covariantes, puisque toutes les fonctions fondamentales
ultérieures deviennent des invariants projectifs de celles-ci. Cela explique
la position distinguée de la forme de courbure riemannienne \([\Omega_2]\)
pour les formes quadratiques. La question de l'équivalence des formes
quadratiques, respectivement des formes de dimension \(p\), se ramène
précisément à l'équivalence de \([\Omega_2]\), respectivement de
\([\Omega_2]\ldots[\Omega_p]\), et de leurs dérivées covariantes, par
transformation linéaire; et l'annulation identique de \([\Omega_2]\),
respectivement de \([\Omega_2]\ldots[\Omega_p]\), est nécessaire et
suffisante pour la transformation en formes à coefficients constants, ce
qui, pour les formes quadratiques, énonce l'un des résultats les plus connus.

% unit: noether.p12.par.017
Une présentation plus détaillée doit paraître prochainement dans les
\emph{Mathematische Annalen}.
